\begin{abstract}
\noindent
Multi-task large language model inference on resource-constrained hardware
faces a fundamental tension: full fine-tuning is prohibitively expensive, yet
existing parameter-efficient serving systems all assume centralised
orchestration: a single datacenter, scheduler, or repository.

This systematic literature review investigates whether a peer-to-peer (P2P)
research direction is theoretically grounded and practically motivated: can
lightweight, task-specific adapters be discovered, retrieved, composed, and
served across autonomous nodes (each hosting a shared frozen backbone) without
any central coordinator? To establish the evidence base, the review follows
PRISMA 2020 and Wohlin (2014) snowballing, surveying 123 papers across five
clusters: PEFT, adapter composition, LLM inference serving, MoE routing, and
P2P and federated learning.

The central finding is that, to the best of the author's knowledge no existing 
system simultaneously addresses a frozen shared backbone, adapter-level exchange, 
P2P topology, decentralised discovery, multi-task fusion, and the absence of a 
central coordinator. The combination of these requirements represents a promising 
but unexplored research direction. The literature is partitioned into four gap 
quadrants: conceptual, algorithmic, systems, and empirical, each identifying 
open research questions that the reviewed literature has not yet addressed 
in an integrated manner.
\end{abstract}

\noindent\textbf{Keywords:} parameter-efficient fine-tuning; low-rank adaptation; adapter composition; 
peer-to-peer systems; decentralised inference; distributed hash table; multi-task learning; 
large language model serving; federated learning; mixture of experts; decentralised adapter exchange; 
capability embeddings; non-IID distribution

\newpage